% !TEX program = xelatex
\documentclass[12pt,a4paper]{article}

\usepackage{fontspec}
\usepackage{polyglossia}
\setmainlanguage{russian}
\setotherlanguage{english}
\setmainfont{DejaVu Serif}
\setsansfont{DejaVu Sans}
\setmonofont{DejaVu Sans Mono}

\usepackage{amsmath,amssymb,amsfonts,mathtools}
\usepackage{bm}
\usepackage{geometry}
\usepackage{enumitem}
\usepackage{graphicx}
\usepackage{hyperref}
\usepackage{microtype}
\geometry{margin=2.5cm}
\hypersetup{colorlinks=true,linkcolor=black,urlcolor=blue,citecolor=black}
\setlength{\emergencystretch}{3em}
\sloppy
\newcommand{\statementbox}[1]{%
\begin{center}%
\fbox{\begin{minipage}{0.88\linewidth}\centering #1\end{minipage}}%
\end{center}%
}

\DeclareMathOperator{\Adm}{Adm}
\DeclareMathOperator{\Struct}{Struct}
\DeclareMathOperator{\Dist}{Dist}
\DeclareMathOperator{\Tr}{Tr}
\DeclareMathOperator{\supp}{supp}
\DeclareMathOperator{\Dom}{Dom}
\DeclareMathOperator{\Cond}{Cond}
\DeclareMathOperator{\Loss}{Loss}
\DeclareMathOperator{\Coh}{Coh}

\newcommand{\TMI}{\mathcal T_{\mathrm{MI}}}
\newcommand{\QGI}{\mathcal QG_I}
\newcommand{\M}{\mathcal M_U}
\newcommand{\MIphys}{\mathcal M_I^{\mathrm{phys}}}
\newcommand{\MIQG}{\mathcal M_I^{\mathrm{QG}}}
\newcommand{\Gspace}{\mathfrak G}
\newcommand{\Hilb}{\mathcal H}
\newcommand{\Prob}{\mathbb P}
\newcommand{\Unc}{\mathcal U}
\newcommand{\Eff}{\mathrm{eff}}
\newcommand{\cl}{\mathrm{cl}}
\newcommand{\Id}{\mathbf 1}
\newcommand{\D}{\mathcal D}
\newcommand{\BorelG}{\mathcal B_G}
\newcommand{\Lag}{\mathcal L}

\title{\textbf{Теория многообразных интерфейсов}\\
\large Осторожный интерфейсный подход к квантовой гравитации}
\author{Рабочая математико-физическая фиксация}
\date{Версия: TMI-QG-1.0 draft}

\begin{document}
\maketitle

\begin{abstract}
Данный текст фиксирует осторожный переход от физического слоя Теории многообразных интерфейсов к квантово-гравитационному слою. Текст не утверждает, что квантовая гравитация уже решена. Его задача скромнее: ввести интерфейсный каркас, в котором метрика перестаёт быть фиксированным фоном, геометрия рассматривается как квантовая возможность, а физический интерфейс описывает переход от неопределённости геометрии к эффективной классической структуре пространства-времени.

Центральная формула:
\[
\boxed{
\Gspace
\Rightarrow
\Psi_G[g]
\Rightarrow
\Prob_G(g)
\Rightarrow
I_G
\Rightarrow
 g_{\Eff}
}
\]

Смысл: пространство возможных геометрий задаёт область геометрической неопределённости; квантовое состояние геометрии задаёт распределение по этой области; гравитационный интерфейс переводит распределение геометрических возможностей в эффективную наблюдаемую геометрию. В версии TMI-QG-1.0 draft дополнительно уточняется линейность линдбладовского генератора, разделяется стандартный марковский режим $\Gamma_I[\chi_I,g;\theta]$ и самосогласованное феноменологическое расширение, добавляется явный расчёт $\chi_I(\tau)\Rightarrow\Gamma_I(\tau)\Rightarrow\eta_I(\tau)$, формулируются основные результаты версии и вводится библиография по связанным формализмам.
\end{abstract}

\tableofcontents

\section{Статус и осторожность формулировки}

Настоящий текст является не заявлением о завершённой квантовой гравитации, а рабочей интерфейсной схемой:
\statementbox{ТМИ-QG --- это интерфейсная модель перехода от квантовой неопределённости геометрии к классической структуре пространства-времени.}

В этой версии используются три уже зафиксированных тезиса физического слоя Теории многообразных интерфейсов:
\[
\boxed{\text{геометрия задаёт область возможного,}}
\]
\[
\boxed{\text{квантовая физика задаёт распределение возможного,}}
\]
\[
\boxed{\text{интерфейс переводит возможность в структуру.}}
\]

Чисто математическая абстракция Теории многообразных интерфейсов вынесена в связанную публикацию:
\begin{quote}
\textit{Mathematical Abstraction of the Theory of Manifold Interfaces}.\\
DOI: \href{https://doi.org/10.5281/zenodo.19978895}{10.5281/zenodo.19978895}.
\end{quote}

Настоящий текст строит физико-математическое расширение этой абстракции в сторону квантовой гравитации.

\section{От фиксированной геометрии к квантовой геометрии}

В классическом физическом слое метрика считается заданной:
\[
g_{\mu\nu}\quad \text{задана на}\quad \M.
\]

Тогда причинная структура, допустимые переходы и квантовое распределение строятся относительно этой метрики:
\[
g
\Rightarrow
\Adm_c(g)
\Rightarrow
\Hilb_I(g)
\Rightarrow
\Psi_I
\Rightarrow
\Dist_{\mathrm{poss}}(\Psi_I;g).
\]

Для перехода к квантовой гравитации ключевое изменение состоит в том, что метрика больше не рассматривается как фиксированный фон:
\[
\boxed{g_{\mu\nu}\ \text{становится квантовой возможностью}.}
\]

Вводится пространство допустимых геометрий:
\[
\boxed{
\Gspace:=\mathcal G_{\mathrm{Lor}}(\M)/\mathrm{Diff}(\M).
}
\]
Здесь $\mathcal G_{\mathrm{Lor}}(\M)$ обозначает класс допустимых лоренцевых геометрий на $\M$.

Факторизация по группе диффеоморфизмов подчёркивает, что физически существенны не координатные представления метрики, а геометрические классы эквивалентности.

\section{Квантовое состояние геометрии}

Квантовое состояние геометрии обозначается как функционал
\[
\Psi_G[g].
\]

Формально можно записать:
\[
\boxed{
|\Psi_G\rangle
=
\int_{\Gspace}
\Psi_G[g]\,|g\rangle\,\D g.
}
\]

Здесь:
\begin{itemize}[leftmargin=2em]
    \item $|g\rangle$ --- формальный базис геометрических альтернатив;
    \item $\Psi_G[g]$ --- амплитуда геометрии $g$;
    \item $\D g$ --- формальная мера на пространстве геометрий.
\end{itemize}

Осторожное замечание: мера $\D g$ не считается в этой версии построенной окончательно. Она обозначает проблему меры в пространстве геометрий. Поэтому запись следует понимать как формальную схему, а не как завершённую конструкцию квантовой гравитации.

Для рабочей версии вводится гильбертово пространство геометрических возможностей:
\[
\boxed{
\Hilb_G
:=
L^2(\Gspace,\BorelG,\mu_G).
}
\]
Здесь $\BorelG$ --- выбранная $\sigma$-алгебра различимых классов геометрий, а $\mu_G$ --- эффективная или регуляризованная мера на $\Gspace$. В строгой версии $\mu_G$ должна быть построена явно; в настоящем тексте она фиксирует место, где проблема меры должна быть решена или ограничена выбором дискретизации, калибровки, регуляризации либо конечного семейства геометрических альтернатив.

Тогда квантовое состояние геометрии можно понимать как
\[
\Psi_G\in\Hilb_G,
\qquad
\|\Psi_G\|^2
=
\int_{\Gspace}|\Psi_G[g]|^2\,d\mu_G(g)=1.
\]
А распределение возможных геометрий задаётся как
\[
\boxed{
d\Prob_G(g)=|\Psi_G[g]|^2\,d\mu_G(g).
}
\]

В дискретизированной или эффективной модели:
\[
\boxed{
|\Psi_G\rangle
=
\sum_i c_i |g_i\rangle,
\qquad
\sum_i |c_i|^2=1.
}
\]

Тогда:
\[
\boxed{
\Prob_G(g_i)=|c_i|^2.
}
\]

\section{Физическое интерфейсное многообразие}

Для перехода от интерпретации к физической модели вводится физическое интерфейсное многообразие:
\[
\boxed{
\MIphys
:=
(\M,g_{\mu\nu},\chi_I,\Adm_I,\Hilb_I,\Prob_{\Psi,I},\Struct_I).
}
\]

В квантово-гравитационном слое одна фиксированная метрика заменяется пространством возможных геометрий:
\[
\boxed{
\MIQG
:=
(\Gspace,\Hilb_G,\Psi_G,\chi_I,\mathcal E_I,\Struct_G).
}
\]

Здесь:
\begin{itemize}[leftmargin=2em]
    \item $\Gspace$ --- пространство допустимых геометрий;
    \item $\Hilb_G=L^2(\Gspace,\BorelG,\mu_G)$ --- рабочее гильбертово пространство геометрических возможностей;
    \item $\Psi_G$ --- квантовое состояние геометрии;
    \item $\chi_I$ --- физическое интерфейсное поле;
    \item $\mathcal E_I$ --- интерфейсный квантовый канал;
    \item $\Struct_G$ --- структурированная геометрия результата.
\end{itemize}

\section{Интерфейсное поле и обратная связь}

Для усиления физического статуса интерфейса вводится не просто функция, а поле как секция интерфейсного расслоения:
\[
\boxed{
E_I\longrightarrow \M,
\qquad
\chi_I\in\Gamma(E_I).
}
\]
Здесь $E_I$ --- интерфейсное расслоение над пространством-временем, а $\Gamma(E_I)$ --- пространство его гладких секций. В минимальной скалярной модели можно взять тривиальное расслоение $E_I=\M\times\mathbb R$, тогда $\chi_I:\M\to\mathbb R$. В более общей версии $\chi_I$ может быть векторным, тензорным или операторнозначным интерфейсным полем.

Локальное значение $\chi_I(x)$ кодирует физическое состояние интерфейсного слоя в точке $x$. Его роль состоит в том, чтобы связывать три уровня: причинную геометрию, квантовое распределение и структуризацию результата.

В классической геометрии причинность задаётся направлением
\[
g_{\mu\nu}\Rightarrow \Adm_c(g).
\]

В интерфейсно-полевой версии добавляется обратная связь:
\[
\boxed{
g_{\mu\nu}\leftrightarrow\chi_I.
}
\]
То есть геометрия ограничивает допустимые интерфейсные переходы, а интерфейсное поле вносит собственный вклад в эффективную геометрию.

Минимальная форма этой обратной связи:
\[
\boxed{
G_{\mu\nu}[g]
=
\frac{8\pi G}{c^4}
\left(
T_{\mu\nu}^{(\Psi)}+T_{\mu\nu}^{(I)}
\right).
}
\]

Здесь $T_{\mu\nu}^{(I)}$ --- эффективный тензор энергии-импульса интерфейсного слоя. В лагранжевой форме он должен выводиться из интерфейсного действия:
\[
\boxed{
T_{\mu\nu}^{(I)}
:=
-\frac{2}{\sqrt{-g}}
\frac{\delta S_I}{\delta g^{\mu\nu}}.
}
\]
Так обратная связь становится не декларацией, а вариационной конструкцией: если интерфейсный слой имеет действие $S_I$, то он вносит вклад в правую часть уравнений геометрии. Если
\[
T_{\mu\nu}^{(I)}\to 0,
\]
то восстанавливается обычная форма уравнений Эйнштейна:
\[
G_{\mu\nu}[g]
=
\frac{8\pi G}{c^4}T_{\mu\nu}^{(\Psi)}.
\]

\section{Квантовый гравитационный интерфейс}

На квантовом уровне интерфейс понимается как переход от неопределённости к структуре. В обычном квантовом слое:
\[
I_Q:\Unc_Q\to\Struct_I.
\]

В квантовой гравитации аналогичная формула принимает вид:
\[
\boxed{
I_G:\Unc_G\to\Struct_G.
}
\]

Где:
\begin{itemize}[leftmargin=2em]
    \item $\Unc_G$ --- геометрическая неопределённость;
    \item $\Struct_G$ --- определённая или эффективная геометрическая структура;
    \item $I_G$ --- гравитационный интерфейс.
\end{itemize}

Для дискретной суперпозиции геометрий:
\[
\boxed{
I_G:
\sum_i c_i |g_i\rangle
\longrightarrow
 g_{\Eff}.
}
\]

Более осторожно:
\[
\boxed{
I_G:
\Dist_{\mathrm{poss}}(\Gspace)
\longrightarrow
\Struct_G.
}
\]

Это означает, что интерфейс не обязан буквально ``разрушать'' суперпозицию. Более корректно:
\statementbox{Интерфейс задаёт механизм структуризации геометрической неопределённости.}

Эта структуризация может проявляться как декогеренция, редукция, эффективный выбор классической геометрии или регистрация геометрического результата относительно конкретной структуры различения.

\section{Запись через матрицу плотности и квантовый канал}

Пусть квантовое состояние геометрии задано матрицей плотности:
\[
\rho_G=|\Psi_G\rangle\langle\Psi_G|.
\]

Интерфейсный переход задаётся квантовым каналом, зависящим от интерфейсного поля:
\[
\boxed{
\mathcal E_I=\mathcal E_{\chi_I},
\qquad
\mathcal E_{\chi_I}(\rho_G)
=
\sum_k K_k[\chi_I,g]\,\rho_G\,K_k[\chi_I,g]^{\dagger}.
}
\]
Условие полной положительности и сохранения следа записывается как
\[
\boxed{
\sum_k K_k[\chi_I,g]^{\dagger}K_k[\chi_I,g]=\Id.
}
\]
Тем самым канал $\mathcal E_I$ не является внешней операцией: он определяется интерфейсной конфигурацией $\chi_I$ и тем, как данная конфигурация различает геометрические альтернативы.

В базовой версии операторы $K_k$ зависят от $\chi_I$ и $g$, но не зависят напрямую от самого состояния $\rho_G$. Это сохраняет линейность квантового канала:
\[
\mathcal E_{\chi_I}(a\rho_1+b\rho_2)
=
a\mathcal E_{\chi_I}(\rho_1)+b\mathcal E_{\chi_I}(\rho_2).
\]
Если в дальнейшем будет использоваться зависимость $K_k$ от $\Psi_G$ или $\rho_G$, такая модель должна рассматриваться отдельно как эффективная нелинейная феноменология, а не как стандартный квантовый канал.

После интерфейсной структуризации:
\[
\boxed{
\rho_G\longmapsto \rho_G'=
\mathcal E_{\chi_I}(\rho_G).
}
\]

В пределе классической геометрии:
\[
\boxed{
\rho_G'
\approx
|g_{\Eff}\rangle\langle g_{\Eff}|.
}
\]

В более реалистичном случае возникает смешанное состояние:
\[
\boxed{
\rho_G'
\approx
\sum_k p_k |g_k\rangle\langle g_k|,
\qquad
p_k=\Tr(E_k^{(I)}\rho_G),
}
\]
где
\[
E_k^{(I)}=K_k[\chi_I,g]^{\dagger}K_k[\chi_I,g].
\]

В этой форме интерфейсный канал описывает не только ``коллапс'', но более общий механизм: декогеренцию, редукцию, coarse-graining или эффективную регистрацию геометрии относительно структуры различения, заданной $\chi_I$.

\section{Интерфейсный функционал действия}

Минимальный классический функционал действия интерфейсно-полевой версии:
\[
\boxed{
S_{\TMI}[g,\Psi,\chi_I]
=
S_{\mathrm{GR}}[g]
+
S_Q[\Psi,g]
+
S_I[\chi_I,g,\Psi].
}
\]

Вариации дают систему уравнений:
\[
\boxed{
\frac{\delta S_{\TMI}}{\delta g^{\mu\nu}}=0,
\qquad
\frac{\delta S_{\TMI}}{\delta \Psi}=0,
\qquad
\frac{\delta S_{\TMI}}{\delta \chi_I}=0.
}
\]

Третье уравнение является новым интерфейсным уравнением:
\[
\boxed{
\frac{\delta S_{\TMI}}{\delta \chi_I}=0.
}
\]

Для минимальной модели интерфейсное действие можно записать как
\[
\boxed{
S_I[\chi_I,g,\Psi]
=
\int_{\M}
\Lag_I(\chi_I,\nabla\chi_I,g,\Psi)\sqrt{-g}\,d^4x.
}
\]
Рабочий скалярный пример:
\[
\boxed{
\Lag_I
=
-\frac{1}{2}g^{\mu\nu}\nabla_\mu\chi_I\nabla_\nu\chi_I
-
V(\chi_I)
-
\chi_I\,J_I(g,\Psi).
}
\]
Здесь $V(\chi_I)$ --- потенциал интерфейсного поля, а $J_I(g,\Psi)$ --- источник или связка интерфейсного поля с геометрическим и квантовым слоями. Эта связка задаёт, как квантовое распределение и геометрическая кривизна влияют на структуризацию.

Вариация по $\chi_I$ даёт минимальное уравнение интерфейсного поля:
\[
\boxed{
\Box_g\chi_I
-
\frac{dV}{d\chi_I}
=
J_I(g,\Psi),
}
\]
где $\Box_g:=\nabla^\mu\nabla_\mu$.

Для скалярной минимальной части тензор энергии-импульса интерфейсного слоя имеет вид
\[
\boxed{
T_{\mu\nu}^{(I)}
=
\nabla_\mu\chi_I\nabla_\nu\chi_I
-
g_{\mu\nu}
\left(
\frac{1}{2}\nabla^\alpha\chi_I\nabla_\alpha\chi_I
+
V(\chi_I)
+
\chi_I J_I(g,\Psi)
\right)
+
T_{\mu\nu}^{(I,\mathrm{coup})}.
}
\]
Последний член $T_{\mu\nu}^{(I,\mathrm{coup})}$ учитывает явную зависимость источника $J_I(g,\Psi)$ от метрики. Если такой зависимости нет или она пренебрежимо мала в эффективной модели, этот член можно опустить.

Для квантово-гравитационной версии вводится формальный интеграл по геометриям и интерфейсным конфигурациям:
\[
\boxed{
\mathcal Z_I
=
\int_{\Gspace}\D g
\int \D\chi_I\,
\exp\left(\frac{i}{\hbar}S_{\TMI}[g,\chi_I,\Psi]\right).
}
\]

Осторожное замечание: $\mathcal Z_I$ является формальным интерфейсным функционалом. В данной версии не утверждается, что мера $\D g$ или интеграл по всем геометриям уже математически завершены.

\section{Главная цепочка ТМИ-QG}

Основная цепочка квантово-гравитационного слоя:
\[
\boxed{
\Gspace
\Rightarrow
\Hilb_G
\Rightarrow
\Psi_G
\Rightarrow
\Dist_{\mathrm{poss}}(\Gspace)
\Rightarrow
I_G
\Rightarrow
\Struct_G
\Rightarrow
 g_{\Eff}.
}
\]

Смысловая расшифровка:
\begin{enumerate}[leftmargin=2em]
    \item $\Gspace$ задаёт область возможных геометрий.
    \item $\Hilb_G$ задаёт пространство квантовых геометрических состояний.
    \item $\Psi_G$ задаёт амплитудную структуру геометрической неопределённости.
    \item $\Dist_{\mathrm{poss}}(\Gspace)$ задаёт распределение возможных геометрий.
    \item $I_G$ задаёт интерфейс структуризации.
    \item $\Struct_G$ является результатом перехода от неопределённости к структуре.
    \item $g_{\Eff}$ является эффективной классической геометрией.
\end{enumerate}

Короткая формула:
\[
\boxed{
\QGI:
\Psi_G[g]
\Rightarrow
\Prob_G(g)
\Rightarrow
I_G
\Rightarrow
 g_{\Eff}.
}
\]

\section{Пределы: возврат к ОТО и квантовой механике}

\subsection{Предел общей теории относительности}

Если интерфейсный вклад исчезает:
\[
T_{\mu\nu}^{(I)}\to0,
\]
и геометрическое состояние становится классически определённым:
\[
\rho_G'\to |g_{\cl}\rangle\langle g_{\cl}|,
\]
то восстанавливается классическая ОТО:
\[
\boxed{
G_{\mu\nu}[g_{\cl}]
=
\frac{8\pi G}{c^4}T_{\mu\nu}.
}
\]

\subsection{Предел квантовой механики на фиксированном фоне}

Если геометрия фиксирована:
\[
\Gspace\to\{g_{\mathrm{fixed}}\},
\]
то квантово-гравитационный слой редуцируется к квантовой теории на заданном фоне:
\[
\boxed{
\Hilb_I(g_{\mathrm{fixed}})
\Rightarrow
\Psi_I
\Rightarrow
\Dist_{\mathrm{poss}}(\Psi_I;g_{\mathrm{fixed}}).
}
\]

В этом пределе ТМИ-QG не конкурирует с обычной квантовой механикой, а содержит её как частный случай при фиксированной геометрии.

\section{Минимальная физическая гипотеза}

Минимальная гипотеза интерфейсного подхода к квантовой гравитации:
\statementbox{Физический гравитационный интерфейс является механизмом перехода от квантовой неопределённости геометрии к эффективной классической структуре пространства-времени.}

Формально:
\[
\boxed{
I_G:
\Unc_G
\longrightarrow
\Struct_G.
}
\]

Или:
\[
\boxed{
\Psi_G[g]
\Rightarrow
\Prob_G(g)
\Rightarrow
I_G
\Rightarrow
 g_{\Eff}.
}
\]


\section{Физический смысл интерфейсного поля}

Важно подчеркнуть: $\chi_I$ не вводится как ещё одно обычное скалярное поле без собственной роли. Его специфическая функция состоит в том, что оно параметризует интерфейсный переход
\[
\boxed{
\begin{gathered}
\text{геометрическая возможность}\\
\longrightarrow\\
\text{квантовое распределение}\\
\longrightarrow\\
\text{структурированный результат}
\end{gathered}
}
\]

Иначе говоря, $\chi_I$ отличается от обычного поля не только возможной динамикой, но и своей операциональной функцией:
\[
\boxed{
\chi_I
\quad\text{задаёт канал структуризации}\quad
\mathcal E_{\chi_I}.
}
\]

Поэтому одна и та же геометрическая суперпозиция может давать разные эффективные структуры при разных интерфейсных конфигурациях:
\[
\boxed{
\rho_G'
=
\mathcal E_{\chi_I}(\rho_G),
\qquad
\rho_G''
=
\mathcal E_{\tilde\chi_I}(\rho_G),
\qquad
\rho_G'\neq \rho_G''.
}
\]

Смысл: интерфейсное поле определяет, какие геометрические альтернативы становятся различимыми, какие когерентности подавляются, какие структуры фиксируются и какая эффективная геометрия возникает после перехода.

В этой версии можно зафиксировать функциональную роль интерфейсного поля так:
\[
\boxed{
\chi_I:
(\Gspace,\Psi_G,\rho_G)
\longmapsto
\mathcal E_{\chi_I}
\longmapsto
\Struct_G.
}
\]

\section{Физически мотивированный источник связи \texorpdfstring{$J_I(g,\Psi)$}{J I(g,Psi)}}

Источник $J_I(g,\Psi)$ должен выражать, какие физические величины управляют интерфейсной структуризацией. В минимальной феноменологической версии можно взять
\[
\boxed{
J_I(g,\Psi)
=
\alpha R
+
\beta\,\langle\widehat T\rangle_\Psi
+
\gamma\,\mathcal C_G^{(I)}.
}
\]

Здесь:
\begin{itemize}[leftmargin=2em]
    \item $R$ --- скалярная кривизна;
    \item $\langle\widehat T\rangle_\Psi$ --- след или скалярная свёртка ожидаемого тензора энергии-импульса квантового слоя;
    \item $\mathcal C_G^{(I)}$ --- интерфейс-зависимая мера когерентности геометрического состояния;
    \item $\alpha,\beta,\gamma$ --- параметры связи.
\end{itemize}

Эта запись не утверждает окончательную форму закона. Она фиксирует минимальную гипотезу:
\statementbox{кривизна, энергия и квантовая когерентность могут управлять скоростью и направлением интерфейсной структуризации геометрии.}

Для меры когерентности можно использовать рабочую дискретную форму. Пусть
\[
\rho_G=\sum_{i,j}\rho_{ij}^{(I)}|g_i\rangle_I\,{}_I\langle g_j|.
\]
Тогда внедиагональная когерентность задаётся как
\[
\boxed{
\mathcal C_G^{(I)}
:=
\sum_{i\neq j}|\rho_{ij}^{(I)}|.
}
\]

Здесь элементы $\rho_{ij}^{(I)}$ берутся в том разложении геометрических альтернатив $\{|g_i\rangle_I\}$, которое задаётся интерфейсом. Поэтому когерентность не является абсолютной величиной:
\[
\boxed{
\mathcal C_G^{(I)}
\quad\text{зависит от интерфейсной структуры различения.}
}
\]
Это не недостаток, а часть смысла ТМИ-QG: интерфейс определяет, какие геометрические альтернативы становятся различимыми и какие когерентности подлежат подавлению.

Параметры $\alpha,\beta,\gamma$ являются размерными константами связи. Они выбираются так, чтобы все слагаемые в $J_I(g,\Psi)$ имели одну и ту же физическую размерность. Поэтому феноменологическая сумма выше должна пониматься не как безразмерное сложение разнородных величин, а как эффективная размерно согласованная модель связи.

В этом случае источник $J_I$ усиливается, когда геометрическое состояние содержит значимую суперпозиционную неопределённость относительно структуры различения, заданной интерфейсом. Это согласуется с идеей, что интерфейсное поле активнее там, где требуется переход от неопределённости к структуре.



\section{Феноменологический статус модели}

Данная версия является эффективной феноменологической моделью, а не завершённой фундаментальной теорией всех геометрических степеней свободы. Это означает, что часть объектов вводится как минимальные рабочие анзацы, позволяющие перейти от общей интерфейсной идеи к вычислимой схеме.

Главная феноменологическая гипотеза состоит в том, что динамика интерфейсного поля $\chi_I$ индуцирует эффективную скорость подавления геометрической когерентности. В стандартном линейном марковском режиме эта скорость записывается как
\[
\boxed{
\Gamma_I=\Gamma_I[\chi_I,g;\theta].
}
\]
Здесь $\theta$ обозначает набор эффективных параметров редуцированного описания: coarse-graining, выбор базиса геометрических альтернатив, масштаб эволюции, параметры волновых пакетов и возможные внешние условия.

В минимальной версии удобно разделить фоновую и интерфейсную части:
\[
\boxed{
\Gamma_I
=
\Gamma_0
+
\delta\Gamma[\chi_I,g;\theta].
}
\]
Здесь $\Gamma_0$ обозначает базовый эффективный уровень декогеренции, который может быть связан с внешней средой, coarse-graining или выбранным приближением, а $\delta\Gamma[\chi_I,g;\theta]$ обозначает собственную интерфейсную поправку.

Для TMI-QG принципиальна именно вторая часть:
\[
\boxed{
\delta\Gamma[\chi_I,g;\theta]
}
\]
поскольку она выражает новую физическую ставку теории: геометрическая структуризация не вводится как внешняя процедура наблюдения, а связывается с физическим интерфейсным слоем.

Следовательно, в этой версии не утверждается, что конкретная формула для $\Gamma_I$ уже выведена из фундаментального действия. Утверждается более осторожное:
\statementbox{существует эффективная интерфейсная поправка $\delta\Gamma[\chi_I,g;\theta]$, которая параметризует скорость перехода от геометрической когерентности к структурированной эффективной геометрии.}

Именно это делает модель проверяемой в принципе: если интерфейсная поправка отлична от нуля,
\[
\boxed{
\delta\Gamma[\chi_I,g;\theta]\neq0,
}
\]
то она должна проявляться как изменение скорости подавления внедиагональных элементов геометрической матрицы плотности.

В стандартном линейном марковском режиме параметры $\theta$ считаются внешними или эффективно заданными параметрами редуцированной модели. Поэтому генератор действует линейно на $\rho_G$. Если в дальнейшем вводится зависимость вида $\Gamma_I[\chi_I,g,\rho_G]$, такая зависимость должна трактоваться отдельно: как самосогласованное нелинейное феноменологическое расширение, а не как стандартный линейный генератор Линдблада. Без этой оговорки такая зависимость могла бы нарушить линейность стандартной открытой квантовой динамики.


\section{Эффективная связь интерфейсного действия \texorpdfstring{$S_I$}{SI} со скоростью структуризации \texorpdfstring{$\Gamma_I$}{Gamma I}}

В версии TMI-QG-1.0 draft связь между интерфейсным действием и скоростью структуризации уточняется как эффективная, а не фундаментально выведенная. Минимальная логика имеет вид
\[
\boxed{
S_I[\chi_I,g,\Psi]
\Rightarrow
\chi_I(\tau)
\Rightarrow
\delta\Gamma[\chi_I,g;\theta]
\Rightarrow
\mathcal E_{\chi_I}.
}
\]

Это означает, что действие $S_I$ задаёт динамику интерфейсного поля, а выбранный эффективный функционал $\Gamma_I[\chi_I,g;\theta]$ переводит эту динамику в параметр квантового канала. Параметры $\theta$ фиксируют coarse-graining, выбор редуцированного пространства, масштаб времени, параметры волновых пакетов и возможные внешние условия. В данной версии не утверждается строгий микроскопический вывод $\Gamma_I$ из $S_I$; фиксируется более осторожный эффективный мост:
\[
\boxed{
\Gamma_I[\chi_I,g;\theta]
=
\Gamma_0+
\delta\Gamma[\chi_I,g;\theta].
}
\]

Фоновая часть $\Gamma_0$ описывает декогеренцию, связанную с coarse-graining, внешней средой или выбором эффективного описания. Интерфейсная часть $\delta\Gamma$ должна исчезать в пределе отсутствия физического интерфейсного вклада:
\[
\boxed{
\chi_I\to\chi_{I,0},
\qquad
\delta\Gamma[\chi_{I,0},g;\theta]\to0.
}
\]

Тем самым обычная квантово-космологическая декогеренция остаётся допустимым пределом, а новая гипотеза TMI-QG состоит только в наличии дополнительной интерфейсной поправки:
\[
\boxed{
\delta\Gamma[\chi_I,g;\theta]\neq0.
}
\]

В минимальной эффективной модели можно считать, что $\delta\Gamma$ является функционалом локальных скаляров, построенных из $\chi_I$, кривизны, интерфейсной структуры различения и эффективных параметров модели:
\[
\boxed{
\delta\Gamma
=
F_{\ge0}
\left(
|\nabla\chi_I|_g^2,
K,
R^2,
\mathcal C_G^{(I)};\theta
\right).
}
\]

Здесь $F_{\ge0}$ задаётся феноменологически. В дальнейшем более строгая версия должна вывести его из микроскопической динамики интерфейсного поля, из интегрирования скрытых степеней свободы или из эффективного влияния интерфейсного слоя на геометрические волновые пакеты.

\section{Эффективный реконструируемый параметр: \texorpdfstring{$\Gamma_I$}{Gamma I}}

Для перехода от каркаса к физической модели нужна величина, которую можно связать не обязательно с прямым измерением, а с реконструкцией по эффектам подавления геометрической когерентности. Поэтому $\Gamma_I$ в версии TMI-QG-1.0 draft не называется фундаментальной наблюдаемой величиной. Она вводится как эффективный реконструируемый параметр интерфейсной структуризации или интерфейсной декогеренции геометрии. Поскольку скорость декогеренции должна быть неотрицательной, минимальный феноменологический анзац имеет вид
\[
\boxed{
\Gamma_I
:=
\lambda_\chi\,\left|g^{\mu\nu}\nabla_\mu\chi_I\nabla_\nu\chi_I\right|
+
\alpha_K |K|
+
\alpha_R R^2
+
\alpha_C\mathcal C_G^{(I)},
\qquad
\Gamma_I\ge0.
}
\]

Эту же запись можно понимать как конкретизацию общей схемы
\[
\boxed{
\Gamma_I
=
\Gamma_0
+
\delta\Gamma[\chi_I,g;\theta]
}
\]
при выборе
\[
\boxed{
\delta\Gamma[\chi_I,g;\theta]
=
\lambda_\chi\left|g^{\mu\nu}\nabla_\mu\chi_I\nabla_\nu\chi_I\right|
+
\alpha_K |K|
+
\alpha_R R^2
+
\alpha_C\mathcal C_G^{(I)}.
}
\]
Если в конкретной модели отсутствует базовая внешняя декогеренция, можно положить $\Gamma_0=0$.

Здесь
\[
K:=R_{\mu\nu\rho\sigma}R^{\mu\nu\rho\sigma}
\]
--- инвариант Кречмана, $R$ --- скалярная кривизна, а $\mathcal C_G^{(I)}$ --- интерфейс-зависимая когерентность. Модуль у градиентного члена введён потому, что в лоренцевой сигнатуре величина
\[
g^{\mu\nu}\nabla_\mu\chi_I\nabla_\nu\chi_I
\]
не обязана быть положительной. В более общей версии можно заменить явную формулу на
\[
\boxed{
\Gamma_I=F_{\ge0}\left(\nabla\chi_I,K,R^2,\mathcal C_G^{(I)}\right),
}
\]
где $F_{\ge0}$ --- неотрицательный феноменологический функционал.

Параметры $\lambda_\chi,\alpha_K,\alpha_R,\alpha_C$ являются размерными константами связи и выбираются так, чтобы $\Gamma_I$ имела размерность обратного времени или обратного собственного параметра эволюции. При положительных коэффициентах правая часть задаёт неотрицательную скорость интерфейсной структуризации.

Физический смысл:
\statementbox{чем сильнее интерфейсная неоднородность, кривизна и интерфейс-зависимая геометрическая когерентность, тем выше может быть скорость перехода от неопределённой геометрии к эффективной структуре.}

Для осторожности $\Gamma_I$ в этой версии следует считать не установленным предсказанием, а эффективным феноменологическим параметром, который потенциально может быть реконструирован из скорости подавления когерентности, ширины волновых пакетов, потери интерференционных членов или других следствий интерфейсной структуризации. Его задача --- показать, как ТМИ-QG может перейти от формальной схемы к проверяемым поправкам, не утверждая пока фундаментальный закон.

Например, если геометрическая матрица плотности подчиняется эффективному уравнению декогеренции,
\[
\boxed{
\frac{d\rho_G}{d\tau}
=
-\frac{i}{\hbar}[\widehat H_G,\rho_G]
-
\Gamma_I\,\mathcal D_G(\rho_G),
}
\]
то $\Gamma_I$ задаёт силу интерфейсного подавления геометрических когерентностей. Здесь $\mathcal D_G$ --- эффективный супероператор декогеренции в пространстве геометрических альтернатив.

\section{Связь скорости структуризации с конкретным квантовым каналом}

Чтобы $\Gamma_I$ не оставалась только феноменологической вставкой, зададим минимальную модель канала чистой декогеренции. Рассмотрим двухуровневое подпространство геометрических альтернатив и параметр
\[
\boxed{
\eta_I(\tau):=e^{-\Gamma_I\tau},
\qquad
0\leq\eta_I\leq1.
}
\]
Величина $\Gamma_I\tau$ должна быть безразмерной; следовательно, если $\tau$ имеет размерность времени или собственного параметра эволюции, то $[\Gamma_I]=[\tau]^{-1}$.

Если скорость интерфейсной структуризации зависит от параметра эволюции, то постоянную экспоненту следует заменить интегральной формой:
\[
\boxed{
\eta_I(\tau)
=
\exp\left(-\int_0^\tau \Gamma_I(s)\,ds\right).
}
\]
Постоянная скорость является частным случаем:
\[
\Gamma_I(s)=\Gamma_I=\mathrm{const}
\quad\Rightarrow\quad
\eta_I(\tau)=e^{-\Gamma_I\tau}.
\]
Именно величина
\[
\int_0^\tau \Gamma_I(s)\,ds
\]
является безразмерной накопленной силой интерфейсной декогеренции.

Минимальный CPTP-канал чистой декогеренции можно записать через операторы Крауса
\[
\boxed{
K_0=
\sqrt{\frac{1+\eta_I}{2}}\,\Id,
\qquad
K_1=
\sqrt{\frac{1-\eta_I}{2}}\,\sigma_z.
}
\]
Тогда
\[
\boxed{
\mathcal E_{\chi_I}(\rho)
=
K_0\rho K_0^\dagger+K_1\rho K_1^\dagger.
}
\]
Так как
\[
K_0^\dagger K_0+K_1^\dagger K_1=\Id,
\]
канал сохраняет след и является физически допустимым квантовым каналом в данной двухсостояний модели.

Если
\[
\rho=
\begin{pmatrix}
\rho_{11} & \rho_{12}\\
\rho_{21} & \rho_{22}
\end{pmatrix},
\]
то действие канала имеет вид
\[
\boxed{
\mathcal E_{\chi_I}(\rho)=
\begin{pmatrix}
\rho_{11} & \eta_I\rho_{12}\\
\eta_I\rho_{21} & \rho_{22}
\end{pmatrix}.
}
\]
Следовательно,
\[
\boxed{
\rho_{12}\longmapsto \eta_I(\tau)\rho_{12},
\qquad
\rho_{21}\longmapsto \eta_I(\tau)\rho_{21}.
}
\]
Именно поэтому $\Gamma_I$ можно интерпретировать как локальную скорость интерфейсного подавления когерентности между геометрическими альтернативами, а $\eta_I(\tau)$ --- как накопленный коэффициент сохранения когерентности. В этой минимальной модели $\Gamma_I$ входит не произвольно, а как параметр конкретного квантового канала, который переводит геометрическую суперпозицию в смешанное состояние относительно интерфейсной структуры различения.



\section{Линдбладовский генератор интерфейсной декогеренции}

Краусова запись канала удобна для конечного шага эволюции. Для динамической формы полезно записать тот же механизм через уравнение Линдблада. В двухуровневом подпространстве геометрических альтернатив минимальный интерфейсный оператор декогеренции можно задать как
\[
\boxed{
L_I[\chi_I,g;\theta]
=
\sqrt{\frac{\Gamma_I[\chi_I,g;\theta]}{2}}\,\sigma_z.
}
\]
В этой стандартной форме $\Gamma_I[\chi_I,g;\theta]$ не зависит функционально от текущей матрицы плотности $\rho_G$; поэтому линдбладовская динамика остаётся линейной по состоянию. Самосогласованные варианты с $\Gamma_I[\chi_I,g,\rho_G]$ можно рассматривать только как отдельные нелинейные феноменологические расширения.

Тогда динамика матрицы плотности имеет вид
\[
\boxed{
\frac{d\rho_G}{d\tau}
=
-
\frac{i}{\hbar}[H_G,\rho_G]
+
L_I\rho_G L_I^\dagger
-
\frac12\{L_I^\dagger L_I,\rho_G\}.
}
\]
Если в данном двухуровневом секторе гамильтониан не создаёт дополнительного смешивания между $|a_1\rangle$ и $|a_2\rangle$, то для внедиагонального элемента получается
\[
\boxed{
\frac{d\rho_{12}}{d\tau}
=
-
\Gamma_I(\tau)\rho_{12}.
}
\]
Отсюда
\[
\boxed{
\rho_{12}(\tau)
=
\exp\left(-\int_0^\tau\Gamma_I(s)ds\right)\rho_{12}(0).
}
\]
То есть линдбладовская форма даёт тот же коэффициент подавления когерентности, что и краусова запись канала:
\[
\boxed{
\eta_I(\tau)
=
\exp\left(-\int_0^\tau\Gamma_I(s)ds\right).
}
\]

В этой записи физический смысл $\Gamma_I$ становится точнее:
\statementbox{$\Gamma_I$ является не просто числом в экспоненте, а коэффициентом генератора открытой квантовой динамики геометрических альтернатив.}

При этом модель сохраняет феноменологический статус: оператор $L_I$ не выводится в данной версии из полной микроскопической теории, а задаёт минимальный эффективный генератор, совместимый с полной положительностью, сохранением следа и интерфейсной интерпретацией структуризации.


\section{Явный расчёт: \texorpdfstring{$\chi_I(\tau)\Rightarrow\Gamma_I(\tau)\Rightarrow\eta_I(\tau)$}{chi to Gamma to eta}}

Чтобы связь интерфейсного поля со скоростью структуризации была не только формальной, введём минимальный одномерный пример. Пусть в редуцированной космологической модели интерфейсное поле зависит только от параметра эволюции $\tau$:
\[
\boxed{
\chi_I(\tau)=\chi_0 e^{-m\tau},
\qquad
m>0.
}
\]
Тогда
\[
\dot\chi_I(\tau)=-m\chi_0e^{-m\tau},
\qquad
\dot\chi_I^2(\tau)=m^2\chi_0^2e^{-2m\tau}.
\]
В минимальной плоской редуцированной модели, где основная интерфейсная поправка задаётся кинетическим членом поля, можно взять
\[
\boxed{
\Gamma_I(\tau)
=
\Gamma_0+
\lambda_\chi m^2\chi_0^2e^{-2m\tau}.
}
\]
Здесь $\Gamma_0\ge0$ --- фоновая эффективная скорость декогеренции, а второе слагаемое --- интерфейсная поправка, исчезающая при затухании поля.

Тогда накопленный коэффициент подавления когерентности равен
\[
\eta_I(\tau)
=
\exp\left(-\int_0^\tau\Gamma_I(s)ds\right).
\]
Подставляя $\Gamma_I(s)$, получаем
\[
\int_0^\tau\Gamma_I(s)ds
=
\Gamma_0\tau
+
\lambda_\chi m^2\chi_0^2
\int_0^\tau e^{-2ms}ds.
\]
Так как
\[
\int_0^\tau e^{-2ms}ds
=
\frac{1-e^{-2m\tau}}{2m},
\]
получаем явную формулу
\[
\boxed{
\eta_I(\tau)
=
\exp\left[
-\Gamma_0\tau
-
\frac{\lambda_\chi m\chi_0^2}{2}
\left(1-e^{-2m\tau}\right)
\right].
}
\]
Следовательно, внедиагональный элемент геометрической матрицы плотности подавляется как
\[
\boxed{
\rho_{12}(\tau)
=
\eta_I(\tau)\rho_{12}(0).
}
\]

Этот пример показывает, как именно поле $\chi_I$ может породить нестационарную скорость интерфейсной структуризации. При $\tau\to\infty$ интерфейсная поправка насыщается, а дальнейшее подавление определяется фоновой частью $\Gamma_0$:
\[
\boxed{
\eta_I(\tau)
\sim
\exp\left(-\Gamma_0\tau-\frac{\lambda_\chi m\chi_0^2}{2}\right).
}
\]
Если $\Gamma_0=0$, то интерфейсное поле даёт конечный накопленный фактор подавления; если $\Gamma_0>0$, когерентность продолжает убывать экспоненциально.

\statementbox{В этом toy-расчёте связь $\chi_I\Rightarrow\Gamma_I\Rightarrow\eta_I$ становится явной: динамика интерфейсного поля задаёт временную скорость подавления геометрической когерентности.}

\begin{figure}[htbp]
    \centering
    \includegraphics[width=0.82\linewidth]{tmi_qg_eta_plot.pdf}
    \caption{Пример накопленного подавления когерентности $\eta_I(\tau)$ для toy-модели $\chi_I(\tau)=\chi_0e^{-m\tau}$. График показывает, как интерфейсная часть быстро даёт конечное подавление, а фоновая скорость $\Gamma_0$ при ненулевом значении продолжает уменьшать когерентность при больших $\tau$.}
    \label{fig:eta-interface}
\end{figure}

\section{Первая вычислимая модель: минисуперпространство}

Чтобы избежать нерешённой проблемы полной меры на всём пространстве геометрий $\Gspace$, можно ввести первую вычислимую модель: минисуперпространственное ограничение. Рассмотрим однородную и изотропную метрику Фридмана--Робертсона--Уокера:
\[
\boxed{
ds^2
=
-N(t)^2dt^2
+
a(t)^2 d\Sigma_k^2.
}
\]

Здесь $a(t)$ --- масштабный фактор, $N(t)$ --- lapse-функция, а $d\Sigma_k^2$ --- метрика трёхмерного пространства постоянной кривизны $k\in\{-1,0,1\}$.

В этой модели пространство геометрий редуцируется:
\[
\boxed{
\Gspace
\longrightarrow
\Gspace_{\mathrm{mini}}
:=
\mathbb R_+,
\qquad
g\longmapsto a.
}
\]

Гильбертово пространство геометрических возможностей становится обычным функциональным пространством:
\[
\boxed{
\Hilb_G^{\mathrm{mini}}
=
L^2(\mathbb R_+,\mu(a)\,da).
}
\]

В минимальной toy-модели можно явно зафиксировать простейший выбор меры:
\[
\boxed{
\mu(a)=1,
\qquad
\Hilb_G^{\mathrm{mini}}=L^2(\mathbb R_+,da).
}
\]
Более общая эффективная мера может быть записана как
\[
\boxed{
\mu(a)=a^p,
\qquad
p\in\mathbb R,
}
\]
где параметр $p$ может кодировать выбранный факторный порядок, редукцию из полной геометрической меры или феноменологическую нормировку. В данной версии, если явно не указано иное, численный toy-пример использует минимальный выбор $\mu(a)=1$.

Квантовое состояние геометрии:
\[
\boxed{
\Psi_G[g]
\longrightarrow
\Psi(a),
\qquad
\int_0^\infty |\Psi(a)|^2\mu(a)\,da=1.
}
\]

Распределение возможных геометрий:
\[
\boxed{
d\Prob_G(a)=|\Psi(a)|^2\mu(a)\,da.
}
\]


В качестве минимального квантово-космологического уравнения можно ввести минисуперпространственную форму уравнения Уилера--ДеВитта:
\[
\boxed{
\widehat{\mathcal H}_{\mathrm{mini}}\Psi(a)=0.
}
\]
Оно выражает квантовое ограничение для геометрической волновой функции $\Psi(a)$. В данной версии не фиксируется единственный факторный порядок оператора $\widehat{\mathcal H}_{\mathrm{mini}}$; достаточно зафиксировать его роль как оператора ограничения, задающего допустимые квантовые состояния геометрии в минисуперпространстве. Схематически можно писать
\[
\boxed{
\widehat{\mathcal H}_{\mathrm{mini}}
=
-\hbar^2 A(a)\frac{d^2}{da^2}
-\hbar^2 B(a)\frac{d}{da}
+U_{\mathrm{mini}}(a),
}
\]
где $A(a)$, $B(a)$ и $U_{\mathrm{mini}}(a)$ зависят от выбранной космологической модели, материи, кривизны $k$ и факторного порядка.


\subsection{Минимальный сценарий: \texorpdfstring{$k=0$, $\Lambda\neq0$}{k=0, Lambda nonzero}}

Для первой физически конкретной версии можно выбрать плоскую FRW-модель с космологической постоянной:
\[
\boxed{
k=0,
\qquad
\Lambda\neq0.
}
\]

В этой редукции потенциал минисуперпространства можно записывать в модельной форме
\[
\boxed{
U_{\mathrm{mini}}(a)
=
U_\Lambda(a)+U_{\mathrm{m}}(a),
}
\]
где $U_\Lambda(a)$ кодирует вклад космологической постоянной, а $U_{\mathrm{m}}(a)$ --- эффективный вклад материи или выбранного часового поля. Без фиксации факторного порядка можно использовать рабочий оператор
\[
\boxed{
\widehat{\mathcal H}_{\mathrm{mini}}
=
-\hbar^2\frac{d^2}{da^2}
+U_{\mathrm{mini}}(a).
}
\]

Тогда допустимые геометрические волновые функции задаются уравнением
\[
\boxed{
\left[-\hbar^2\frac{d^2}{da^2}+U_{\mathrm{mini}}(a)\right]\Psi(a)=0.
}
\]

Эта запись не претендует на единственный правильный факторный порядок и не заменяет полную квантовую космологию. Её задача --- дать минимальную площадку, где можно явно связать волновые пакеты $\psi_1(a),\psi_2(a)$, интерфейсное поле $\chi_I(t)$ и скорость подавления когерентности $\Gamma_I$.


Интерфейсное поле в самой простой космологической версии также редуцируется к функции времени:
\[
\boxed{
\chi_I(x)
\longrightarrow
\chi_I(t).
}
\]

Тогда интерфейсный переход принимает вычислимую форму:
\[
\boxed{
\Psi(a)
\Rightarrow
\Prob_G(a)
\Rightarrow
\mathcal E_{\chi_I(t)}
\Rightarrow
 a_{\Eff}.
}
\]

Эффективный масштабный фактор можно задать как среднее по структурированному состоянию:
\[
\boxed{
a_{\Eff}
:=
\Tr(\rho_G'\,\widehat a),
\qquad
\rho_G'=\mathcal E_{\chi_I}(\rho_G).
}
\]

В дискретной модели:
\[
\boxed{
a_{\Eff}
=
\sum_i p_i' a_i,
\qquad
p_i'=\langle a_i|\rho_G'|a_i\rangle.
}
\]

Таким образом, первая вычислимая версия ТМИ-QG имеет вид:
\[
\boxed{
\begin{gathered}
\text{суперпозиция масштабных факторов}\\
\longrightarrow\\
\text{интерфейсная структуризация}\\
\longrightarrow\\
\text{эффективная космологическая геометрия}
\end{gathered}
}
\]

Эта модель не решает полную квантовую гравитацию, но создаёт минимальную площадку для расчётов: можно исследовать, как $\chi_I(t)$, $J_I(a,\Psi)$ и $\Gamma_I(a,\chi_I)$ влияют на переход от $\Psi(a)$ к $a_{\Eff}$.



\subsection{Связь двухсостояний модели с волновой функцией \texorpdfstring{$\Psi(a)$}{Psi(a)}}

Двухсостояний пример не следует понимать как произвольную замену непрерывной волновой функции. Его можно получить как эффективное приближение, если волновая функция масштабного фактора имеет две хорошо различимые локализованные компоненты:
\[
\boxed{
\Psi(a)=c_1\psi_1(a)+c_2\psi_2(a),
\qquad
\|\psi_1\|=\|\psi_2\|=1.
}
\]
Здесь $\psi_1(a)$ сосредоточена в окрестности $a_1$, а $\psi_2(a)$ --- в окрестности $a_2$:
\[
\boxed{
\supp(\psi_1)\approx a_1,
\qquad
\supp(\psi_2)\approx a_2,
\qquad
\langle\psi_1|\psi_2\rangle\approx0.
}
\]
Тогда подпространство
\[
\boxed{
\Hilb_2:=\mathrm{span}\{\psi_1,\psi_2\}
}
\]
можно отождествить с двухуровневой моделью геометрических альтернатив:
\[
\boxed{
|a_1\rangle\equiv\psi_1(a),
\qquad
|a_2\rangle\equiv\psi_2(a).
}
\]
В этом смысле дискретная модель является не отдельной онтологией, а эффективной coarse-grained редукцией непрерывного состояния $\Psi(a)$ к двум различимым геометрическим пакетам.

\subsection{Двухсостояний пример}

Для первого явного расчётного примера ограничим минисуперпространство двумя геометрическими альтернативами:
\[
|a_1\rangle,
\qquad
|a_2\rangle.
\]
Пусть начальное состояние масштабного фактора имеет вид
\[
\boxed{
|\Psi\rangle
=
c_1|a_1\rangle+c_2|a_2\rangle,
\qquad
|c_1|^2+|c_2|^2=1.
}
\]
Тогда начальная матрица плотности:
\[
\rho_G
=
\begin{pmatrix}
|c_1|^2 & c_1\overline{c_2}\\
c_2\overline{c_1} & |c_2|^2
\end{pmatrix}.
\]
Применяя канал чистой декогеренции из предыдущего раздела, получаем
\[
\boxed{
\rho_G'(\tau)
=
\begin{pmatrix}
|c_1|^2 & \eta_I(\tau)c_1\overline{c_2}\\
\eta_I(\tau)c_2\overline{c_1} & |c_2|^2
\end{pmatrix}.
}
\]
При $\eta_I(\tau)\to0$ когерентность между $|a_1\rangle$ и $|a_2\rangle$ эффективно подавляется. При постоянной скорости это соответствует $\Gamma_I\tau\gg1$:
\[
\rho_G'(\tau)
\longrightarrow
\begin{pmatrix}
|c_1|^2 & 0\\
0 & |c_2|^2
\end{pmatrix}.
\]

Важно различать два объекта.

Первый объект --- зарегистрированный результат:
\[
\boxed{
a_{\mathrm{reg}}\in\{a_1,a_2\},
\qquad
\Prob(a_{\mathrm{reg}}=a_k)=p_k'=
\langle a_k|\rho_G'|a_k\rangle.
}
\]
Он соответствует одной реализованной геометрической альтернативе в акте регистрации.

Второй объект --- эффективное среднее, или coarse-grained классическая величина:
\[
\boxed{
a_{\Eff}
=
\Tr(\rho_G'\widehat a)
=
p_1'a_1+p_2'a_2.
}
\]
Если канал чистой декогеренции не меняет диагональные вероятности, то
\[
\boxed{
a_{\Eff}
=
|c_1|^2a_1+|c_2|^2a_2.
}
\]
Следовательно, $a_{\mathrm{reg}}$ обозначает отдельную зарегистрированную альтернативу, а $a_{\Eff}$ обозначает эффективное среднее описание после подавления когерентности. Смешивать эти два уровня нельзя: один относится к событию регистрации, другой --- к эффективной классической геометрии.


\subsection{Численный toy-пример}

Возьмём минимальную калибровку через относительную флуктуацию масштабного фактора:
\[
\boxed{
a_1=1,
\qquad
\delta a=1,
\qquad
a_2=a_1+\delta a=2,
\qquad
|c_1|^2=0.7,
\qquad
|c_2|^2=0.3.
}
\]
Здесь $a_1$ задаёт выбранную единицу масштаба, а $\delta a$ --- относительное отличие второй геометрической альтернативы. Тогда
\[
\boxed{
a_{\Eff}=0.7\cdot a_1+0.3\cdot(a_1+\delta a)=a_1+0.3\delta a=1.3.
}
\]
В общем двухсостояний случае
\[
\boxed{
a_{\Eff}=a_1+|c_2|^2\delta a.
}
\]
Если для простоты взять $c_1=\sqrt{0.7}$ и $c_2=\sqrt{0.3}$, то начальная величина внедиагонального элемента равна
\[
\boxed{
|\rho_{12}(0)|=\sqrt{0.7\cdot0.3}\approx0.458.
}
\]
После интерфейсной декогеренции:
\[
\boxed{
|\rho_{12}(\tau)|=\eta_I(\tau)\sqrt{0.21}.
}
\]
Для постоянной скорости $\Gamma_I=\mathrm{const}$ получаем $\eta_I=e^{-\Gamma_I\tau}$. Например:
\[
\begin{array}{c|c|c}
\Gamma_I\tau & e^{-\Gamma_I\tau} & |\rho_{12}(\tau)| \\
\hline
0 & 1 & 0.458 \\
1 & 0.368 & 0.169 \\
3 & 0.050 & 0.023 \\
5 & 0.007 & 0.003
\end{array}
\]
Этот toy-пример показывает физический смысл параметра $\Gamma_I$: при росте накопленной величины $\int_0^\tau\Gamma_I(s)ds$ геометрическая когерентность быстро подавляется, но диагональные вероятности $0.7$ и $0.3$ сохраняются. Поэтому интерфейсная структуризация не равна произвольному выбору результата; она задаёт переход от суперпозиционной когерентности к статистически интерпретируемой смеси геометрических альтернатив.

\section{Что именно является новой гипотезой ТМИ-QG}

Важно отделить стандартные математические элементы от новой гипотезы ТМИ-QG. Квантовые каналы, декогеренция, матрицы плотности, минисуперпространство и уравнение Уилера--ДеВитта уже существуют в известных формализмах квантовой теории, открытых квантовых систем и квантовой космологии. Поэтому новизна данной версии не состоит в том, что заново вводится декогеренция или заново формулируется минисуперпространство.

Новая гипотеза состоит в связке:
\[
\boxed{
\chi_I
\Rightarrow
\Gamma_I
\Rightarrow
\mathcal E_{\chi_I}
\Rightarrow
\rho_G'
\Rightarrow
 g_{\Eff}.
}
\]
То есть физическое интерфейсное поле $\chi_I$ рассматривается как структурирующий слой, который задаёт скорость подавления геометрической когерентности, параметризует квантовый канал и тем самым участвует в переходе от геометрической неопределённости к эффективной классической геометрии.

В самой короткой форме новая гипотеза ТМИ-QG записывается так:
\statementbox{Интерфейсное поле не просто сопровождает геометрию, а параметризует структуризацию квантовой геометрии.}

Поэтому стандартные элементы выполняют роль математического языка, а новая физическая ставка ТМИ-QG заключается в том, что канал структуризации геометрии имеет интерфейсную природу:
\statementbox{Канал $\mathcal E_{\chi_I}$ является не внешней процедурой наблюдения, а эффективным проявлением физического интерфейсного слоя.}

Эта формулировка сохраняет осторожность: TMI-QG-1.0 draft не объявляет завершённую квантовую гравитацию, а задаёт проверяемую исследовательскую программу. Её дальнейшая задача --- вывести или ограничить $\Gamma_I$ из динамики $\chi_I$, а затем сравнить предсказания модели с известными пределами квантовой космологии, полуклассической гравитации и теории декогеренции.




\section{Что модель не утверждает}

Для публикационной ясности важно отделить положительную гипотезу TMI-QG от тех утверждений, которые данная версия не делает. Версия TMI-QG-1.0 draft не должна читаться как завершённое решение всех проблем квантовой гравитации. Её статус --- минимальная феноменологическая и вычислимая модель интерфейсной структуризации квантовой геометрии.

В частности, данная версия не утверждает:
\begin{enumerate}[leftmargin=2em]
    \item что построена полная мера $\mu_G$ на всём пространстве геометрий $\Gspace$;
    \item что решена проблема времени в квантовой гравитации;
    \item что декогеренция сама по себе решает проблему выбора единственного результата;
    \item что $\Gamma_I$ уже фундаментально выведена из микроскопического действия;
    \item что минисуперпространственная модель исчерпывает все геометрические степени свободы;
    \item что TMI-QG заменяет существующие подходы к квантовой гравитации.
\end{enumerate}

Положительное утверждение версии значительно уже:
\statementbox{TMI-QG-1.0 draft предлагает эффективный интерфейсно-полевой механизм, в котором $\chi_I$ параметризует скорость геометрической декогеренции, а канал $\mathcal E_{\chi_I}$ описывает переход от квантовой геометрической когерентности к эффективной структуре.}

\section{Связь с существующими физико-математическими формализмами}

TMI-QG-1.0 draft использует несколько известных математических языков как инструменты, но не отождествляется с ними полностью.

\begin{enumerate}[leftmargin=2em]
    \item \textbf{Уравнение Уилера--ДеВитта и квантовая космология~\cite{DeWitt1967,Wheeler1968,Kiefer2012}.} Минисуперпространственная часть документа использует стандартную идею квантового ограничения
    \[
    \widehat{\mathcal H}_{\mathrm{mini}}\Psi(a)=0,
    \]
    но новая гипотеза состоит не в самом уравнении, а в добавлении интерфейсного механизма структуризации геометрических альтернатив.

    \item \textbf{Открытые квантовые системы~\cite{Lindblad1976,Kraus1983}.} Краусовы операторы и линдбладовский генератор используются как стандартный язык описания полностью положительной динамики и декогеренции. Новая часть TMI-QG --- интерпретация параметров канала как функций интерфейсного поля:
    \[
    \chi_I\Rightarrow\Gamma_I\Rightarrow\mathcal E_{\chi_I}.
    \]

    \item \textbf{Декогеренция~\cite{Zurek2003,Joos2003}.} Модель не утверждает, что декогеренция впервые вводится в квантовой гравитации. Она утверждает более узко: геометрическая декогеренция может быть параметризована интерфейсным полем, которое задаёт структуру различения между геометрическими альтернативами.

    \item \textbf{Полуклассическая гравитация~\cite{BirrellDavies1982}.} В пределе подавленной геометрической когерентности возникает эффективная геометрия $g_{\Eff}$, что позволяет рассматривать обычные уравнения для классической метрики как предельное описание.

    \item \textbf{Суммирование по геометриям.} Формальная запись через $\int_{\Gspace}\D g$ близка по духу к суммированию по геометриям, но TMI-QG делает акцент не на полной мере, а на интерфейсном переходе от распределения геометрий к эффективной структуре.
\end{enumerate}

Итоговая позиция:
\statementbox{TMI-QG-1.0 draft не заменяет существующие подходы, а предлагает интерфейсный слой, который связывает квантовую геометрическую неопределённость, декогеренцию и появление эффективной классической структуры.}

\section{Пределы применимости}

Данная версия имеет ограниченную область применимости. Она не описывает полное пространство всех геометрических степеней свободы и не заменяет завершённую теорию квантовой гравитации. Её естественный статус --- эффективная феноменологическая схема в минисуперпространстве и в конечномерных coarse-grained подпространствах геометрических альтернатив.

Основные пределы применимости:
\begin{enumerate}[leftmargin=2em]
    \item Модель надёжнее всего работает в редуцированной области
    \[
    \boxed{
    \Gspace_{\mathrm{mini}}=\mathbb R_+,
    }
    \]
    где геометрическая степень свободы представлена масштабным фактором $a$.

    \item Двухсостояний пример следует понимать как эффективное приближение двух хорошо различимых волновых пакетов, а не как утверждение, что пространство-время фундаментально двухсостоянийно.

    \item Параметр $\Gamma_I$ является эффективным реконструируемым параметром. Он может быть связан с наблюдаемыми следствиями только через подавление когерентности, параметры канала, ширины волновых пакетов или другие физические эффекты.

    \item Формула для $\Gamma_I$ является минимальным анзацем. Она должна быть либо выведена из более фундаментального действия, либо ограничена экспериментальными, космологическими или математическими условиями.

    \item Полная мера $\mu_G$ на пространстве геометрий $\Gspace$ не строится в этой версии. Вместо неё используется минимальная мера минисуперпространства $\mu(a)=1$ или параметрическое семейство $\mu(a)=a^p$.
\end{enumerate}

Поэтому наиболее точная формулировка статуса версии:
\statementbox{TMI-QG-1.0 draft является минимальной феноменологической моделью интерфейсной структуризации квантовой геометрии в редуцированном пространстве геометрических альтернатив с явным линдбладовским генератором и конкретным минисуперпространственным сценарием.}


\section{Основные результаты версии TMI-QG-1.0 draft}

Версия TMI-QG-1.0 draft фиксирует следующие результаты:
\begin{enumerate}[leftmargin=2em]
    \item Квантовая геометрия описывается как распределение по пространству геометрических альтернатив, редуцированному в минимальной модели к минисуперпространству $\mathbb R_+$.
    \item Физическое интерфейсное поле $\chi_I$ задаёт не саму геометрию напрямую, а эффективную скорость структуризации геометрической когерентности.
    \item В стандартном линейном режиме скорость декогеренции записывается как $\Gamma_I[\chi_I,g;\theta]$, где $\theta$ фиксирует параметры редуцированного описания.
    \item Канал структуризации может быть задан как CPTP-канал чистой декогеренции или как линдбладовский генератор с оператором $L_I=\sqrt{\Gamma_I/2}\,\sigma_z$.
    \item Для внедиагональных элементов геометрической матрицы плотности получается закон
    \[
    \rho_{12}(\tau)
    =
    \exp\left(-\int_0^\tau\Gamma_I(s)ds\right)\rho_{12}(0).
    \]
    \item Toy-расчёт с $\chi_I(\tau)=\chi_0e^{-m\tau}$ показывает явную цепочку
    \[
    \chi_I(\tau)
    \Rightarrow
    \Gamma_I(\tau)
    \Rightarrow
    \eta_I(\tau)
    \Rightarrow
    \rho_G'(\tau).
    \]
    \item Новая гипотеза TMI-QG состоит не в повторном введении декогеренции, а в том, что геометрическая декогеренция параметризуется физическим интерфейсным слоем.
\end{enumerate}

Итоговая формула версии:
\[
\boxed{
S_I
\Rightarrow
\chi_I(\tau)
\Rightarrow
\Gamma_I(\tau)
\Rightarrow
\mathcal E_{\chi_I}
\Rightarrow
\rho_G'
\Rightarrow
 g_{\Eff}.
}
\]

\section{Открытые задачи}

Чтобы данная схема стала полноценной физической теорией, необходимо решить следующие задачи:
\begin{enumerate}[leftmargin=2em]
    \item Построить или явно ограничить меру $\mu_G$ на полном пространстве геометрий $\Gspace$.
    \item Исследовать минимальные меры минисуперпространства $\mu(a)=1$ и $\mu(a)=a^p$, а также связать параметр $p$ с факторным порядком или эффективной редукцией полной геометрической меры.
    \item Развить минисуперпространственную модель $\Gspace_{\mathrm{mini}}=\mathbb R_+$ за пределы двухсостояний приближения.
    \item Связать двухсостояний пример с конкретными решениями или приближёнными волновыми пакетами $\psi_1(a),\psi_2(a)$ уравнения Уилера--ДеВитта.
    \item Вывести параметры канала $\mathcal E_{\chi_I}$ и эффективный параметр $\Gamma_I$ из более фундаментальной динамики поля $\chi_I$.
    \item Уточнить разложение $\Gamma_I=\Gamma_0+\delta\Gamma[\chi_I,g;\theta]$ и физический смысл фоновой части $\Gamma_0$.
    \item Исследовать динамическую форму $\eta_I(\tau)=\exp(-\int_0^\tau\Gamma_I(s)ds)$ для нестационарных интерфейсных режимов.
    \item Уточнить физическую размерность $\chi_I$, $J_I$, $\Gamma_I$ и всех констант связи.
    \item Проверить предельные случаи:
    \[
    T_{\mu\nu}^{(I)}\to0,
    \qquad
    \Gamma_I\to0,
    \qquad
    \Gamma_I\to\infty.
    \]
    \item Сравнить модель с известными подходами: полуклассической гравитацией, квантовой космологией, декогеренцией, петлевой квантовой гравитацией и суммированием по геометриям.
\end{enumerate}

\section{Итог}

Квантовая гравитация в осторожной интерфейсной форме ТМИ определяется как теория перехода:
\statementbox{Геометрическая неопределённость $\longrightarrow$ распределение возможных геометрий\\[0.2em]
$\longrightarrow$ интерфейсная структуризация $\longrightarrow$ эффективное пространство-время.}

Главная формула:
\[
\boxed{
\Gspace
\Rightarrow
\Psi_G[g]
\Rightarrow
\Prob_G(g)
\Rightarrow
I_G
\Rightarrow
 g_{\Eff}.
}
\]

Коротко:
\statementbox{Квантовая гравитация в ТМИ --- это интерфейсный переход от неопределённости геометрии к структуре пространства-времени.}

Версия TMI-QG-1.0 draft делает следующий шаг: добавляет эффективную связь $S_I\to\Gamma_I$, вводит линдбладовский генератор интерфейсной декогеренции, конкретизирует минимальный минисуперпространственный сценарий $k=0$, $\Lambda\neq0$, и явно связывает модель с существующими формализмами квантовой космологии, открытых квантовых систем и декогеренции. Итоговая вычислимая цепочка принимает вид
\[
\boxed{
\Psi(a)
\Rightarrow
\Prob_G(a)
\Rightarrow
\mathcal E_{\chi_I(t)}
\Rightarrow
\rho_G'
\Rightarrow
\{a_{\mathrm{reg}},a_{\Eff}\}.
}
\]


\begin{thebibliography}{99}

\bibitem{DeWitt1967}
B.~S.~DeWitt.
\newblock Quantum theory of gravity. I. The canonical theory.
\newblock \emph{Physical Review}, 160, 1113--1148, 1967.

\bibitem{Wheeler1968}
J.~A.~Wheeler.
\newblock Superspace and the nature of quantum geometrodynamics.
\newblock In C.~M.~DeWitt and J.~A.~Wheeler, editors, \emph{Battelle Rencontres: 1967 Lectures in Mathematics and Physics}. Benjamin, 1968.

\bibitem{Kiefer2012}
C.~Kiefer.
\newblock \emph{Quantum Gravity}.
\newblock Oxford University Press, 3rd edition, 2012.

\bibitem{Lindblad1976}
G.~Lindblad.
\newblock On the generators of quantum dynamical semigroups.
\newblock \emph{Communications in Mathematical Physics}, 48, 119--130, 1976.

\bibitem{Kraus1983}
K.~Kraus.
\newblock \emph{States, Effects, and Operations: Fundamental Notions of Quantum Theory}.
\newblock Springer, 1983.

\bibitem{Zurek2003}
W.~H.~Zurek.
\newblock Decoherence, einselection, and the quantum origins of the classical.
\newblock \emph{Reviews of Modern Physics}, 75, 715--775, 2003.

\bibitem{Joos2003}
E.~Joos, H.~D.~Zeh, C.~Kiefer, D.~Giulini, J.~Kupsch, and I.-O.~Stamatescu.
\newblock \emph{Decoherence and the Appearance of a Classical World in Quantum Theory}.
\newblock Springer, 2nd edition, 2003.

\bibitem{BirrellDavies1982}
N.~D.~Birrell and P.~C.~W.~Davies.
\newblock \emph{Quantum Fields in Curved Space}.
\newblock Cambridge University Press, 1982.

\bibitem{HartleHawking1983}
J.~B.~Hartle and S.~W.~Hawking.
\newblock Wave function of the Universe.
\newblock \emph{Physical Review D}, 28, 2960--2975, 1983.

\end{thebibliography}

\end{document}
